\section{Introduction}
A robot hand that holds an object also hides it. Fingers cover part of the surface, and whatever faces the palm or points away from the camera cannot be seen at all. For most learned in-hand controllers this is a nuisance for state estimation and nothing more, because the goal is given in advance as a target pose or a goal image~\cite{bhardwaj2026viserdex,bhardwaj2026goaldex,kedia2026simtoolreal}. It becomes the central problem once the goal is given in language. Asked to ``present the face marked T,'' the robot may be looking at a cube whose T is on the far side, and nothing in the current image tells it which orientation the instruction refers to (Fig.~\ref{fig:teaser}).

The robot has to look before it can act: turn the object, take another view, read the mark, and only then commit to an orientation. This is active perception~\cite{bajcsy1988active,bohg2017interactive}, except that the viewpoint is moved by contact-rich finger motion and a poor choice can drop the object. A vision-language model can reason about which view to obtain next, but it takes seconds to answer (about five in our trials), and during those seconds the object still has to be tracked and held at every control step, under exactly the finger occlusion that inspection produces.

REVEAL treats this as two visual processes at different rates. A slow \emph{semantic} process decides what is on the object: an agent reads the instruction and the current image, chooses the direction and duration of each inspection, and, once a view shows the requested mark, links that surface to a motor goal. We call this \emph{target binding}, and it is what turns ``the face marked T'' into a skill the hand can execute. A fast \emph{geometric} process tracks where the object is: an RGB-D network fuses image features with measured surface points to estimate pose at every control step, and this pose drives a goal-conditioned policy for the Revo3 hand. Two teacher-supervised local controllers bridge the two rates, holding the object stable while a reply is pending and correcting the policy's commands once a goal is bound.

Controlled studies change one ingredient at a time, and each of these parts earns its place: fresh images, not repeated ones, account for the benefit of feedback; agent-chosen inspections beat a fixed sweep; learned slowing with joint correction beats both fixed slowing and uninterrupted execution. As a complete system, REVEAL reaches \QOneChangedSuccess\% mean success on six simulated objects under changed target-to-skill mappings, \QOneChangedGain{} points above the strongest baseline, and \HardwareAimiSuccess\% over 80 real-hand trials on four objects, where a baseline carrying the same local controllers reaches only \HardwareFineTunedCtrlSuccess\%. Our contributions are:
\begin{itemize}
\item An agentic framework for in-hand visual exploration, in which the agent chooses how to look and binds a visual target to a motor goal from successive observations.
\item A perception and control layer that holds the object steady while the agent reasons, combining an occlusion-robust RGB-D pose estimator with teacher-supervised local controllers.
\item Controlled studies isolating fresh observations, inspection choice and waiting control, with six-object, changed-mapping, two-language and four-object real-hand evaluations.
\end{itemize}
